\paragraph{経路の変更と定数}
経路のデータは各辺の剰余行列の組の関数である。
節点での基底変更に対しては辺の行列式指標を掛けて変換し、
縦辺の列が分裂単射でないところでは零とする。
曲がりの外にある行列を固定し、曲がりの二辺のデータを
可逆補完の末行と先頭列で移す。補完を変えると列フレームが
右から基底変更され、その行列式による重みの変化が
対応式の \(\chi(\det g)\) を打ち消す。従って移した重みは
補完に依らず、次の経路でも同じ変換則を満たす。
対角上では \(Y^{j+1}_j=\widetilde Y^{j+1}_j\) 上にシーソーで周期を移す。
余分の直線の核を随伴に直すと、その列は有限 Fourier 変換を受ける。
これで両経路の二辺を \eqref{eq:exchange-square} の形にそろえ、
交換の後で再びシーソーを使って周期を元の形に戻せる。
係数の項数は一つの剰余列の全個数 \(N_\mathrm{col}\) 以下だけ増える。
実際、補完による重みは単位根であり、直線の Fourier 変換でも
一列について加法指標の単位根を掛けて足すだけである。
\(m^2\) 回の曲がりで増える項数は \(N_\mathrm{col}^{m^2}\) 以下で、
各曲がりの分母に入る非零の交換係数 \(c_{ij}\) を掛け合わせたものが
\(c_{\mathrm{path}}\) である。
\leanlinks{\leansrc{New/Determinant/AuxiliaryForms/Turns}{可逆補完で移した経路データの変換則}
  \leanlinksep \leansrc{New/Determinant/AuxiliaryForms/DiagonalTurns}{対角上の直線と \(m^2\) 回の曲がり}}

\paragraph{終経路の和の分母と大きさ}
\(n_m=\dim X_m\)、\(d_b=b r^4\) とする。終経路のどの辺の
Young 図形も箱の総数が \(d_b\) 以下であり、階数は \(2n_m\) 以下である。
各辺の有限データは一つの剰余行列の指示関数である。
指示関数を有限 Fourier 展開で標準データに戻す分母は、
\(\pi_M^b\) の下の行列の剰余類の個数と、その共役側の同じ個数を掛ければ払える。
捻りの分岐素点に由来する因子は、許容する \(\ell\nmid N_\xi\) では単元である。
さらに微分一箱あたり \(M^b d\) を掛ければ整真空の微分の分母を払える。
従って各辺の核に共通の整数
\[
 D_\mathrm{edge}=(M^b)^{2f(2n_m)^2}(M^b d)^{d_b}
\]
を掛ければ、整な値を整な値へ送る。
終点の二空間の階数は一つだけ違う。二つの最高重みが交錯するので、
各列の最高重みベクトルを分岐の小行列式に展開し、
整数係数の分岐行列を作る。これを双対の Weyl 加群の対と合成したものが、
終経路の二つの値を対にする行列である。
最後の列の基底には、法 \(\mathfrak m=\pi_M^bQ\) を生成するイデール
\(t_{\mathfrak m}\) と真空の \(\alpha_\mathrm{W}\) からなる
\(u=\overline{t_{\mathfrak m}\alpha_\mathrm{W}}/(t_{\mathfrak m}\alpha_\mathrm{W})\)
の \(m\) 乗が入る。\(t_{\mathfrak m}\) は \(M^b N_\xi\) を割り、
\(q_1\alpha_\mathrm{W}\) と \(\alpha_\mathrm{W}^{-1}\) は整なので、
\(M^bq_1\) を掛けた \(u,u^{-1}\) は \(\ell\nmid N_\xi\) で整である。
行列式指標の指数の絶対値を \(e_\tau^\mathrm{last}\) とすると、
この変数の分母は \(\prod_\tau(M^b q_1)^{m e_\tau^\mathrm{last}}\) で払える。
類和の安定化群の位数の分母も含めた一つの整数を \(D_\mathrm{end}\) と取る。
各定数の対数を足すと \(\log D_\mathrm{end}\le x^{98}b^2\)。
両側で \(m\) 回ずつ核を適用するので、
\[
 D_\mathrm{path}=D_\mathrm{end}D_\mathrm{edge}^{2m}
       \le \exp(x^{99}b^2)
\]
を終経路の和に掛ければ整になる。経路データの係数は単位根の和なので
ここでは分母を増やさない。
\leanlinks{\leansrc{New/Determinant/AuxiliaryForms/LastPath}{交錯する重みの整数分岐行列}
  \leanlinksep \leansrc{New/Determinant/AuxiliaryForms/SpanningForms}{終経路の分母を払う整数}}

無捻りの場合に \(G=\exp(x^{99}b^2)\) と置く。
各辺では、箱数 \(d_b\) と階数 \(2n_m\) を
第 \ref{sec:upper-bound} 節の微分ノルムの評価に入れると、
共役でのノルムの倍率は \(G\) 以下である。出発点の二乗ノルムの上界を
\(B_-,B_+\) とすると、対を作る直前のノルムは
\(G^m\sqrt{B_-},G^m\sqrt{B_+}\) 以下になる。
分岐行列のノルムの上界は \((d_b^{d_b})^3\) である。
両経路の \(2m+1\) 個の節点での類数はそれぞれ \(G\) 以下であり、
最後の重みに現れる単位根の個数は
\(N_\mathrm{col}^{m^2}N_\mathrm{mat}^{2m}\) 以下である。
\(N_\mathrm{mat}\) は辺の剰余行列の個数の上界である。
各項の上界を足して
\[
 |\sigma(\Xi)|\le
 N_\mathrm{col}^{m^2}N_\mathrm{mat}^{2m}(d_b^{d_b})^3G^{4m+1}\sqrt{B_-B_+}
 \le\exp(x^{200}b^2)\sqrt{B_-B_+}.
\]
最後は許容データでの \(m,f,r,\log M\) の上界を使う。
双対デルタ形式でテストする側には \(\sqrt{B_-}\) として
その値のノルムが入る。有限側の \(D_\mathrm{path}\) と双対のフック数は、
無限素点では合わせても
\(\exp(x^{201}b^2-x^{200}b^2)\) 以下である。
そこで双対ノルムで割れば、終経路の和から元の値への
ノルムの倍率は \(\exp(x^{201}b^2)\) 以下になる。
深いレベルの安定化群は自明であるため、類ごとのデルタ形式で空間を張る。
短い格子基底を Weyl 加群に作用させ、単位根での多項式の値を選んで
各類の値を張ると、整で共役の二乗ノルムが \(G\) 以下の生成系を得る。
これを全射 \(s\) に適用して、補助形式に必要な二つの評価が得られる。
\leanlinks{\leansrc{New/Determinant/AuxiliaryForms/SpanningForms}{類数、項数および双対評価}
  \leanlinksep \leansrc{New/UpperBound/Size/ShortColumns}{短いデルタ形式による生成}}
